% Adapted from count-t.tex to return values and construct a dictionary.
\begin{blocksection}
\question Write \lstinline{count_t}, which returns the number of lowercase
\lstinline{"t"} characters in \lstinline{word}. Then use it to construct a
dictionary mapping each word in \lstinline{words} to its count.

\begin{lstlisting}
def count_t(word: str) -> int:
    """
    >>> count_t("tatter")
    3
    >>> count_t("tree")
    1
    >>> count_t("")
    0
    """
    _______________________________

    for ____________________________:

        if ____________________________:

            __________________________________

    _______________________________
\end{lstlisting}
\begin{solution}
\begin{lstlisting}
def count_t(word: str) -> int:
    count = 0
    for c in word:
        if c == 't':
            count += 1
    return count
\end{lstlisting}
\end{solution}
\end{blocksection}

\begin{blocksection}
Complete the dictionary comprehension. Its result should be
\lstinline|{'tatter': 3, 'tree': 1}|.

\begin{lstlisting}
words = ["tatter", "tree"]
counts: dict[str, int] = {
    ____________: ____________ for word in words
}
\end{lstlisting}
\begin{solution}
\begin{lstlisting}
counts: dict[str, int] = {
    word: count_t(word) for word in words
}
\end{lstlisting}
\end{solution}
\end{blocksection}

\begin{questionmeta}
\textbf{Teaching Tips}\par\nopagebreak[4]
Suggested Time: 7 min; Difficulty: Easy
\nopagebreak[4]
\begin{itemize}
    \item \textbf{Read the types:} The function takes a string and returns
    an integer. The dictionary has string keys and integer values.
    \item \textbf{Plan the loop:} Start at zero, traverse the characters,
    count each lowercase \lstinline{'t'}, and return the total.
    \item \textbf{Construct the dictionary:} Identify the
    \lstinline{key: value} pair and the iterable. Each word is a key;
    its count is the associated value. The expression produces a dictionary.
    \item \textbf{Connect the two stages:} Ask what \lstinline{count_t(word)}
    returns for one word before evaluating the whole comprehension.
    Updating the local integer count does not alter the input string.
\end{itemize}
\end{questionmeta}
